\documentclass[11pt]{article}

\usepackage[T1]{fontenc}
\usepackage[utf8]{inputenc}
\usepackage{lmodern}
\usepackage[english]{babel}
\usepackage{amsthm,amsmath,amssymb,amsfonts,mathtools,mathrsfs}
\usepackage{geometry}
\geometry{verbose,tmargin=2.1cm,bmargin=2.1cm,lmargin=2.3cm,rmargin=2.3cm,marginparwidth=1.5cm}
\usepackage[shortlabels]{enumitem}
\usepackage[babel]{microtype}
\usepackage{xurl,needspace}
\usepackage[hidelinks]{hyperref}
\usepackage[capitalise]{cleveref}
\usepackage{bookmark}

\setlength{\emergencystretch}{2em}

% Theorem environments.
\newtheorem{theorem}{Theorem}[section]
\newtheorem{proposition}[theorem]{Proposition}
\newtheorem{lemma}[theorem]{Lemma}
\newtheorem{corollary}[theorem]{Corollary}
\newtheorem{claim}[theorem]{Claim}
\theoremstyle{definition}
\newtheorem{defn}[theorem]{Definition}
\newtheorem{inputthm}[theorem]{Proposition}
\theoremstyle{remark}
\newtheorem{remark}[theorem]{Remark}
\crefname{inputthm}{Proposition}{Propositions}
\Crefname{inputthm}{Proposition}{Propositions}

% Notation.
\newcommand{\F}{\mathcal F_q}
\newcommand{\FF}{\mathbb F}
\newcommand{\HH}{\mathrm H}
\newcommand{\PP}{\mathbb P}
\newcommand{\CC}{\mathcal C}
\newcommand{\BB}{\mathcal B}
\newcommand{\DD}{\mathcal D}
\newcommand{\SSS}{\mathscr S}
\newcommand{\eps}{\varepsilon}
\newcommand{\src}[1]{\ifmmode\text{\texttt{\detokenize{#1}}}\else\texttt{\detokenize{#1}}\fi}
\newcommand{\repo}{https://github.com/openai/math/blob/adc7f1241b42e322a6451854ab7e4b4c146bf78a/lean/OAI/Combinatorics/SharpRamsey/}
\newcommand{\oaifile}[2]{\href{\repo#1}{\src{#2}}}

\hypersetup{
  pdftitle={The Ramsey number r(4,4,t) via single-colour compression},
  pdfauthor={},
  pdfkeywords={multicolour Ramsey numbers, multicolor Ramsey numbers, off-diagonal Ramsey theory, three-colour Ramsey numbers, r(4,4,t), K4-free graphs, common independent sets, flag graphs, finite geometry, entropy method, single-colour compression}
}

\title{The Ramsey number $r(4,4,t)$ via single-colour compression}
\author{}
\date{}

\begin{document}
\maketitle
\vspace{-2em}

\begin{abstract}
We prove that, for every $\eps>0$, there exists $c_\eps>0$ such that
$r(4,4,t)\ge c_\eps t^5/(\log t)^{4+\eps}$ for all sufficiently large
integers $t$. Together with the upper bound of He and Wigderson, this
gives $r(4,4,t)=t^5/(\log t)^{4+o(1)}$. We use two flag labellings to
construct two $K_4$-free graphs on a common vertex set. If every pair of
labellings admits a large common independent set, then the two label
sequences selected on that set have large joint entropy. The
single-colour scanning inequality gives a nonnegative excess cost in
which the entropy difference is allowed to be negative. Retaining its
sign and summing over the two colours provides both the budget for
single-colour compression and the final contradiction. The argument
tracks only the total cost of finite candidate tables and alternately
replaces the table for each colour. We use the single-colour scanning
and compression results as inputs, prove the two-colour argument in the
main text, and detail the source interfaces in an appendix.

\medskip
\noindent\textbf{Keywords:} Multicolour Ramsey numbers; off-diagonal
Ramsey theory; three-colour Ramsey numbers; $r(4,4,t)$; $K_4$-free
graphs; common independent sets; flag graphs; finite geometry;
entropy method; single-colour compression
\end{abstract}

\section{Introduction}\label{sec:introduction}

For a positive integer $t$, the Ramsey number $r(4,4,t)$ is the smallest
integer $N$ such that every three-colouring of the edges of $K_N$
contains a $K_4$ in the first colour, a $K_4$ in the second colour, or a
$K_t$ in the third colour. A lower bound can therefore be obtained by
constructing two $K_4$-free graphs $G_1,G_2$ on the same vertex set such
that $G_1\cup G_2$ has no large independent set. The graphs are allowed
to share edges: when forming the three-colouring, all common edges may
be assigned to the first colour.

Our construction uses two copies of the flag labelling underlying
Brada\v{c}'s directed graph $D^*(3,q)$~\cite{Bradac}. The geometric
construction ensures that each of the resulting graphs is $K_4$-free.
To control their common independent sets, we use the single-colour
scanning and compression results in the OpenAI
source~\cite{OAI-scan,OAI-step}. Our main result is the following.

\begin{theorem}\label{thm:main}
As $t\to\infty$,
\[
 r(4,4,t)=\frac{t^5}{(\log t)^{4+o(1)}}.
\]
\end{theorem}

The upper bound follows from the unconditional result of He and
Wigderson~\cite[Theorem~4 and Appendix~B]{HW}, which gives
$r(4,4,t)\le Ct^5/(\log t)^4$ for an absolute constant $C$.

We briefly describe the proof. Suppose that every pair of flag
labellings has a common independent set of the prescribed size. Select
one such set in each configuration, always listing its vertices in
increasing order. The two resulting label sequences are supported on
exactly the same positions. A counting argument then gives a lower
bound on their joint entropy in which the cost of selecting these
positions is paid only once.

For each colour, we retain a family of candidate tables specifying the
possible label at each position. The single-colour scanning inequality
associates a nonnegative excess cost with each such family. Crucially,
the entropy difference occurring in this cost need not itself be
nonnegative. By keeping its sign, summing the two excess costs, and
using subadditivity of entropy, we obtain one inequality that both
bounds the cost available to each compression step and yields a
contradiction once the total table cost is small. A local decomposition
separates sequences concentrated in a few incidence blocks from those
to which the compression input applies. We then replace the table of
one colour at a time, while the other colour records only the selected
subsequence positions. Finitely many rounds, with the number of rounds
depending only on the fixed error parameter, reduce the total table
cost to the required scale.

The rest of the paper is organised as follows. In
Section~\ref{sec:flags}, we describe the flag construction and fix the
parameters. Section~\ref{sec:single} states the single-colour input.
In Section~\ref{sec:balance}, we derive the joint-entropy inequality,
and in Section~\ref{sec:replacement}, we prove the single-colour
replacement lemma. Section~\ref{sec:main-proof} completes the proof of
Theorem~\ref{thm:main}. Appendix~\ref{app:source} details the source
interfaces and parameter substitutions for the single-colour input.
Appendix~\ref{app:other} records the scope of several related results
and explains why they are not used as substitutes for that input.

\paragraph{Status of the manuscript.}
The two-colour extension presented here remains a research draft
awaiting independent verification. The sources and parameter
substitutions for the single-colour input are recorded in
Appendix~\ref{app:source}, while the new two-colour argument is given in
the main text. We have neither rebuilt the full OpenAI Lean dependency
tree nor formalised the new argument.

\section{Flag labels and the graph construction}\label{sec:flags}

Let $q$ be a prime, and let $V=\FF_q^4$. We call the one-dimensional
subspaces of $V$ \emph{points} and its three-dimensional subspaces
\emph{hyperplanes}. A \emph{flag} is a pair $(p,H)$ with $p\subset H$,
and we denote the set of all flags by $\F$. Define a directed graph on
$\F$ by
\begin{equation}\label{eq:arc}
 (p,H)\longrightarrow(p',H')
 \quad\Longleftrightarrow\quad
 p'\subset H\ \text{ and }\ p\not\subset H'.
\end{equation}
With respect to the standard dot product, the identification
$(a,b)=(H^\perp,p)$ gives Brada\v{c}'s graph $D^*(3,q)$. By
\cite[Lemma~2.11, with parameter~$3$]{Bradac}, this directed graph
contains no transitive tournament on four vertices.

Given a flag $f_v$ at each vertex $v\in[N]$, form a graph by joining
$u<v$ precisely when $f_u\to f_v$. This graph is $K_4$-free. Indeed,
the labels of a four-vertex clique, listed in vertex order, would form
a transitive tournament on four vertices. Different vertices may
receive the same flag, but there is no arc between equal flags, so the
labels of a clique are necessarily distinct.

If an independent set is listed in increasing vertex order, then its
flag sequence $((p_i,H_i))_{i=1}^n$ satisfies
\begin{equation}\label{eq:legal}
 i<j,\quad p_j\subset H_i
 \quad\Longrightarrow\quad p_i\subset H_j.
\end{equation}
We call a sequence satisfying \eqref{eq:legal} \emph{legal}. Every
order-preserving subsequence of a legal sequence is again legal. The
direction in \eqref{eq:legal} is exactly the condition
\src{ScanConsistent} in the single-colour input; no reversal of the
sequence is required.

For a subspace $W\le V$, put
\[
 R(W)=\{(p,H):p\subseteq W\subseteq H\}.
\]
We shall use the following geometric estimates:
\begin{equation}\label{eq:geometry}
 \begin{gathered}
 v:=|\F|=(q^3+q^2+q+1)(q^2+q+1),
 \qquad q^5\le v\le4q^5,\\
 |R(W)|\le4q^2,
 \qquad |\{W:W\le V\}|\le q^{16}.
 \end{gathered}
\end{equation}
The first line counts the flags. The estimates for the sets $R(W)$
also appear in~\cite{OAI-geometry}; the number of subspaces can be
bounded by counting row spaces of $4\times4$ matrices. We call $R(W)$
an \emph{incidence block}. There is no arc between any two flags in an
incidence block. Thus labels lying in a small union of blocks can be
described by recording the list of blocks. These are precisely the
orthogonal rectangles in the original compression interface.
Throughout the paper, the occupancy of a block is counted
\emph{with multiplicity in the sequence positions}: repeated labels
are counted once for each occurrence.

Fix $0<\eta\le1$, and set
\begin{equation}\label{eq:scale}
 \begin{gathered}
 \sigma=\log q,\qquad Q=q\sigma^{1+\eta},\qquad
 k=\lfloor Q\rfloor,\\
 N=\lfloor q^5\sigma\rfloor,\qquad
 \beta=\eta/10^7.
 \end{gathered}
\end{equation}
Fix a sufficiently large absolute constant $P\ge64$ that dominates the
fixed prefactor in the original compression theorem, and then set
$c=10^4+2\log P$. The parameter $\beta$ is taken from the single-colour
input and is used only to leave gaps between the logarithmic powers.
All logarithms are natural. From this point onwards, we first fix
$\eta$ and the relevant constants, and then take $q$ sufficiently
large. In particular, we may assume that $\sigma\ge1$ and
$Q/2\le k\le Q$.

The parameter $N$ always denotes the number of vertices of the original
graph, and $k$ is the size of the independent set to be excluded. The
letter $n$ denotes the current sequence length, which may decrease
during the replacement procedure.

\section{The single-colour input}\label{sec:single}

A \emph{candidate table} specifies the allowed flags at each position.
More precisely, let $\Omega$ be a nonempty finite set, and suppose that
each $\omega\in\Omega$ determines a legal sequence $F(\omega)$ of length
$n\ge1$. Let $u(\omega)\in\CC$ be its table index, where $\CC$ is a
finite family, and assume that
\[
 F(\omega)_j\in S(u(\omega),j),\qquad
 |S(u,j)|\le Pq^3e^\Delta,\qquad \Delta\ge0.
\]
The index $u(\omega)$ may depend on the entire configuration $\omega$.
Once the table family is fixed, however, each entry $S(u,j)$ depends
only on the table index $u$ and the position $j$. A table restricts the
possible labels but need not determine the sequence uniquely. Unused
entries may be replaced by singleton sets.

Equip $\Omega$ with the uniform distribution, and let $\HH(F)$ denote
the Shannon entropy of the \emph{entire sequence}. For a finite random
variable $Z$, we use
\[
 \HH(Z)=-\sum_z\PP(Z=z)\log\PP(Z=z),
\]
with the convention $0\log0=0$. Define
\begin{equation}\label{eq:defect}
 \DD(F,\CC,\Delta)
 =\log|\CC|+cn+cq\sigma\Delta+3n\sigma-\HH(F).
\end{equation}
The term $\log|\CC|$ accounts for selecting a table, and $\Delta$
measures the increase in its entry sizes relative to $q^3$. The entropy
difference $3n\sigma-\HH(F)$ is allowed to be negative. The input below
asserts the nonnegativity of the \emph{whole} quantity $\DD$, not of
this difference alone. Retaining its sign is essential when we sum the
costs for the two colours.

\Needspace{7\baselineskip}
\begin{inputthm}[Single-colour scanning and finite-table compression]
\label{input:one}
The data above satisfy the following statements.

\begin{enumerate}[label=\textup{(\roman*)},leftmargin=2em]
\item\label{item:nonnegative}
If $\sigma\ge1$, then $\DD(F,\CC,\Delta)\ge0$. No upper bound on the
occupancy of an incidence block is required.

\item\label{item:compression}
Fix $a>0$ and $B_0\ge1$. For all sufficiently large primes $q$, suppose
that
\[
 aQ\le n\le k,\qquad
 \sigma^\beta\le D\le\sigma^{1-\eta/2},
\]
that every sequence occupies at most $B_0\sigma$ positions in each
incidence block $R(W)$, and that
\begin{equation}\label{eq:input-budget}
 \DD(F,\CC,\Delta)\le B_0nD\sigma^{-\beta}.
\end{equation}
Then there exist a subset $\Omega'\subseteq\Omega$, a fixed length $n'$,
and a new table family $\mathcal M$ such that
\begin{equation}\label{eq:output}
 \begin{gathered}
 |\Omega'|\ge|\Omega|/192,\qquad n'\ge n/6144,\\
 \log|\mathcal M|\le2D\sigma^{-\eta/3}n+2,\qquad
 |A(m,j)|\le Pq^3\exp(16D\sigma^{6\beta}).
 \end{gathered}
\end{equation}
For every $\omega\in\Omega'$, the new sequence is an order-preserving
subsequence of $F(\omega)$, and its $j$th entry belongs to
$A(m(\omega),j)$. The length $n'$ is the same for all retained
configurations, but the selected positions may depend on $\omega$.
The new candidate set is determined by $m$ and $j$ alone; the old table
index is not needed. This does not assert probabilistic independence
between the old and new table indices. The threshold for $q$ is uniform
over all the finite sets, tables, sequences $F$, and parameters
$\Delta,D$ satisfying these assumptions.
\end{enumerate}
\end{inputthm}

Part~\textup{(i)} follows from the single-colour scanning inequality
in~\cite{OAI-scan}. Part~\textup{(ii)} follows from the single-colour
compression theorem in~\cite{OAI-step}, together with a conversion to
a finite family of tables. The precise interfaces and parameter
substitutions are given in Appendix~\ref{app:source}. The two parts
serve different purposes: part~\textup{(i)} gives a nonnegative excess
cost for every legal sequence, whereas part~\textup{(ii)} compresses
only sequences with small block occupancies. We therefore perform a
local decomposition before applying part~\textup{(ii)}.

\section{A joint-entropy inequality}\label{sec:balance}

Let $\SSS=(\F^2)^{[N]}$ be the set of all original configurations.
Each configuration assigns two flags to every vertex, one for each
colour. Let $\varnothing\ne\Omega\subseteq\SSS$. Suppose that, for every
$\omega\in\Omega$, we have selected positions
$v_1(\omega)<\cdots<v_n(\omega)$ such that the two label sequences
$X_1,X_2$ on \emph{these same positions} are legal. The choice of
positions may depend on the whole configuration, but the positions
must not be chosen separately for the two colours.

For colour $i$, retain a table family $\CC_i$ and an entry-size
parameter $\Delta_i$, and write
\begin{equation}\label{eq:table-budget}
 \begin{gathered}
 B_i=\log|\CC_i|+cq\sigma\Delta_i,\qquad
 \BB=B_1+B_2,\\
 d_i=B_i+cn+3n\sigma-\HH(X_i).
 \end{gathered}
\end{equation}
Here $B_i$ is the cost of the table family, and $d_i$ is the nonnegative
excess cost in Proposition~\ref{input:one}. Across stages, we track
$B_i$, not the previous value of $d_i$. Whenever configurations are
discarded, we equip the new set $\Omega$ with its uniform distribution
and recompute the entropies. Restricting the set of configurations does
not increase the number of available tables, but it may change these
entropies.

\Needspace{14\baselineskip}
\begin{lemma}\label{lem:balance}
Fix $a,\delta\in(0,1]$. If
\[
 aQ\le n\le k,\qquad |\Omega|\ge\delta|\SSS|,
\]
then
\begin{equation}\label{eq:balance}
 0\le d_1+d_2
 \le\BB-\eta n\log\sigma+C_{a,\delta}n,
\end{equation}
where one may take
$C_{a,\delta}=2c+1+\log(1/a)+\log(1/\delta)$.
In particular, each $d_i\le\BB+C_{a,\delta}n$, giving an input budget
for the next compression. On the other hand, if $a,\delta$ are fixed
and $q\to\infty$, then $\BB=O(n)$ is impossible.
\end{lemma}

\begin{proof}
Any fixed sequence of flag pairs is obtained from at most
$\binom Nn v^{2(N-n)}$ original configurations: choose the common
positions, and then assign arbitrary labels to all remaining vertices.
Since the two colours use the same positions, we pay for only
\emph{one} factor $\binom Nn$. The increasing order is already fixed,
so no additional factor $n!$ is needed. Consequently, under the uniform
distribution on $\Omega$, the probability of any value of $(X_1,X_2)$
is at most $\delta^{-1}\binom Nn v^{-2n}$.

Entropy is at least the negative logarithm of the largest atom.
Therefore,
\begin{align}
 \HH(X_1,X_2)
 &\ge2n\log v-\log\binom Nn+\log\delta\notag\\
 &\ge6n\sigma+\eta n\log\sigma
       +n(\log a-1)+\log\delta.
 \label{eq:joint}
\end{align}
For the second inequality, we used $v\ge q^5$, $N\le q^5\sigma$,
$n\ge aQ$, and $\binom Nn\le(eN/n)^n$.

By Proposition~\ref{input:one}\textup{(i)}, both $d_1$ and $d_2$ are
nonnegative. No independence assumption on the two colours is needed:
subadditivity of entropy~\cite[Theorem~2.6.6]{CT} gives
\begin{align*}
 d_1+d_2
 &=\BB+2cn+6n\sigma-\HH(X_1)-\HH(X_2)\\
 &\le\BB+2cn+6n\sigma-\HH(X_1,X_2).
\end{align*}
Substituting \eqref{eq:joint} proves \eqref{eq:balance}. The remaining
assertions follow from $d_i\ge0$ and $\log\sigma\to\infty$.
\end{proof}

Inequality~\eqref{eq:balance} is the link between the two-colour
construction and the single-colour input. Compression requires a bound
on each individual $d_i$, whereas the configuration count gives a bound
on the joint entropy. Summing the two nonnegative excess costs allows
us to use the joint-entropy estimate to control both of them.

\section{Replacing one colour at a time}\label{sec:replacement}

We first describe how to restrict an existing table to a subsequence.
Suppose that the retained positions are given by a strictly increasing
map $I:[n']\hookrightarrow[n]$. Including $I$ in the old table index
gives the restricted table
\begin{equation}\label{eq:restrict-table}
 \begin{gathered}
 S^{\mathrm{res}}((u,I),j)=S(u,I(j)),\\
 \log|\CC^{\mathrm{res}}|
 \le\log|\CC|+\log\binom n{n'}
 \le\log|\CC|+n\log2.
 \end{gathered}
\end{equation}
The entry-size parameter $\Delta$ does not change. We record positions
\emph{within the current sequence}; we do not reorder the $N$ original
vertices.

\Needspace{15\baselineskip}
\begin{lemma}[Single-colour replacement]\label{lem:rewrite}
Let $\rho=1/24576$, $\theta=1/384$, and
$\lambda=\sigma^{-\eta/4}$. Fix $a,\delta\in(0,1]$ and $K\ge1$.
Suppose that
\[
 aQ\le n\le k,\qquad |\Omega|\ge\delta|\SSS|,\qquad
 \BB\le Kkb,
\]
where $b\ge1$ and $D=\sigma^\beta b\le\sigma^{1-\eta/2}$.
For all sufficiently large $q$, the table of either prescribed colour
$i\in\{1,2\}$ can be replaced so that
\begin{equation}\label{eq:rewrite}
 \begin{gathered}
 |\Omega'|\ge\theta|\Omega|,\qquad n'\ge\rho n,\\
 B_i'\le20k+\tfrac14\lambda kb,\qquad
 B_{3-i}'\le B_{3-i}+k.
 \end{gathered}
\end{equation}
The two colours retain the same original vertex positions throughout.
The constant $K$ affects only the threshold for $q$.
\end{lemma}

\begin{proof}
We prove the statement for $i=1$. We first split each current sequence
into a part concentrated in a small collection of incidence blocks,
which can be described directly by listing the blocks, and a part with
small occupancy in every block, to which the compression input applies.

Let $h=\lceil\sigma\rceil$. For each configuration, repeatedly record
an incidence block containing at least $h$ remaining positions, and
remove all remaining positions in that block. Fix an order on the
blocks so that the procedure is deterministic. Each step removes at
least $h$ new positions, so at most $r=\lfloor n/h\rfloor$ blocks are
recorded. Pad the list to length $r$ with an empty symbol. By
\eqref{eq:geometry},
\begin{equation}\label{eq:peel}
 \begin{gathered}
 \log(\text{number of block lists})
 \le r\log(q^{16}+1)\le17k,\\
 \left|\bigcup\text{recorded blocks}\right|
 \le4q^2n/h\le4q^3\sigma^\eta.
 \end{gathered}
\end{equation}
The positions not removed occupy fewer than $h\le2\sigma$ positions
in each incidence block.

At least one of the removed and remaining parts contains
$n_0=\lfloor n/2\rfloor$ positions. Choose a branch that is available
for at least half the configurations, and let $\Omega_0$ be the set of
those configurations. In each configuration, retain the first $n_0$
positions in the chosen part, simultaneously for both colours. Denote
the resulting sequences by $X_{1,0},X_{2,0}$ and the position selection
by $J$. Then $|\Omega_0|\ge|\Omega|/2$ and $n_0\ge n/4$. The branch and
the length are fixed, but $J$ may still depend on the configuration.

\medskip
\noindent\textbf{Case 1: the concentrated branch.}
Suppose that we retain positions from the removed part. Use the block
list as the new table index for the first colour, with every table
entry equal to the union of the listed blocks. Taking
$\Delta_1'=\eta\log\sigma$, we obtain from \eqref{eq:peel} that
\[
 B_1'\le17k+c\eta q\sigma\log\sigma\le18k
\]
for sufficiently large $q$, since $q\sigma/k\le2\sigma^{-\eta}$.
Thus this branch already reduces the first-colour cost to $O(k)$,
without invoking the compression input.

\medskip
\noindent\textbf{Case 2: the dispersed branch.}
Suppose that we retain positions from the part that was not removed.
Temporarily include $J$ in both table indices as in
\eqref{eq:restrict-table}, and denote the resulting table families by
$\CC_{1,0},\CC_{2,0}$. Their total cost is at most
$\BB+2n\log2$. Apply Lemma~\ref{lem:balance} \emph{afresh} under the
uniform distribution on $\Omega_0$, with parameters $a/4$ and
$\delta/2$. This gives
\[
 \DD(X_{1,0},\CC_{1,0},\Delta_1)
 \le\BB+O_{a,\delta,P}(k)
 \le C_{a,\delta,P,K}kb.
\]
Write $C_*$ for the last constant. Since $k/n_0\le4/a$ and
$D\sigma^{-\beta}=b$, choosing a fixed
$B_0\ge\max\{2,4C_*/a\}$ ensures \eqref{eq:input-budget} and also
covers the block-occupancy bound $2\sigma$. We may therefore apply
Proposition~\ref{input:one}\textup{(ii)}.

The new first-colour table satisfies
\begin{align*}
 \frac{B_1'}k
 &\le\frac2k+2b\sigma^{\beta-\eta/3}
             +32cb\sigma^{7\beta-\eta}\\
 &\le1+\tfrac14\lambda b.
\end{align*}
Here the new entry-size parameter is
$\Delta_1'=16D\sigma^{6\beta}$. The final inequality follows from
$\eta/3-\beta>\eta/4$ and $\eta-7\beta>\eta/4$, once $q$ is sufficiently
large. We retain at least a $1/(2\cdot192)$ fraction of the
configurations and at least a $1/(4\cdot6144)$ fraction of the original
sequence length.

\medskip
In either case, the new candidate sets depend only on the new table
index and the new position, so the old first-colour table can be
discarded immediately. Let $I:[n']\hookrightarrow[n]$ denote the final
position selection for the entire operation. For the second colour,
it is enough to retain $(u_2,I)$ as in \eqref{eq:restrict-table}. Hence
\[
 B_2'\le B_2+\log\binom n{n'}\le B_2+k.
\]
The map $I$ is the composition of the preliminary selection $J$ and
the selection inside the compression step, so these need not be
recorded separately. The original vertex indices remain attached to
each configuration. They need not be included in the table index,
since \eqref{eq:joint} already accounts for the selection of original
positions. This proves \eqref{eq:rewrite}.
\end{proof}

\section{Proof of the main theorem}\label{sec:main-proof}

We now iterate Lemma~\ref{lem:rewrite} a bounded number of times and
obtain the required pair of graphs.

\begin{proposition}\label{prop:construction}
For every fixed $0<\eta\le1$ and all sufficiently large primes $q$,
there exist two $K_4$-free graphs $G_1,G_2$ on the same $N$-vertex set
such that $\alpha(G_1\cup G_2)<k$, where $N$ and $k$ are defined in
\eqref{eq:scale}.
\end{proposition}

\begin{proof}
Each original configuration produces two graphs, one from each
coordinate, using \eqref{eq:arc}. Suppose, for a contradiction, that
no required configuration exists. Thus every configuration has a
common independent set of size $k$. Choose one according to a fixed
rule and list its vertices in increasing order, obtaining two legal
label sequences.

Initially, $\Omega=\SSS$, and each colour has a single table whose
entries are all of $\F$. We may take
\[
 \Delta_1=\Delta_2=2\sigma,\qquad
 \BB_0=4cq\sigma^2,\qquad
 b_*=1+\BB_0/k\le1+8c\sigma^{1-\eta}.
\]

In each round, replace the first-colour table and then the second-colour
table. We always normalise costs by the original parameter $k$, not by
the decreasing current length. Suppose that a round starts with
$\BB\le kb$. For the first replacement, take $K=1$ in
Lemma~\ref{lem:rewrite}. The intermediate total cost is at most
$B_2+21k+\lambda kb/4\le30kb$. For the second replacement, take $K=30$
and use the same $b$ and $D$. This is permitted because changing $K$
only changes the fixed budget constant and the threshold for $q$, not
the output bound of the lemma. When the second colour is replaced, the
first colour pays at most one further position-selection cost. Thus
\begin{equation}\label{eq:round}
 \BB'\le41k+\tfrac12\lambda kb\le k(50+\lambda b).
\end{equation}

To make all thresholds uniform, first fix
\[
 T=\lceil4/\eta\rceil,\qquad
 a=\rho^{2T}/4,\qquad \delta=\theta^{2T},
\]
and only then choose $q$ sufficiently large. There are at most $2T$
single-colour replacements, so throughout the process,
\[
 n\ge\rho^{2T}k\ge aQ,\qquad
 |\Omega|\ge\theta^{2T}|\SSS|=\delta|\SSS|.
\]
The number of rounds depends only on the fixed parameter $\eta$.
Although the retained proportions may be small, they do not decrease
with $q$.

Prescribe
\[
 b_j=100+\lambda^j b_*\quad(0\le j\le T),\qquad
 D_j=\sigma^\beta b_j\quad(0\le j<T).
\]
When $\lambda\le1/2$, we have $50+\lambda b_j\le b_{j+1}$. It follows
inductively from \eqref{eq:round} that the total cost after $j$ rounds
is at most $kb_j$. All compression parameters are admissible, since
\[
 \sigma^\beta\le D_j\le\sigma^\beta(100+b_*)
 =O_P(\sigma^{1-\eta+\beta})
 =o(\sigma^{1-\eta/2}).
\]

After $T$ rounds, the inequality $T\eta/4\ge1$ gives
\[
 b_T\le100+\sigma^{-1}(1+8c\sigma^{1-\eta})=O_P(1).
\]
Consequently, $\BB_T=O_P(k)$. Since the final length satisfies
$n\ge\rho^{2T}k$, this is also $\BB_T=O_\eta(n)$. Applying
Lemma~\ref{lem:balance} at the final stage yields
\[
 0\le d_1+d_2\le-\eta n\log\sigma+O_\eta(n),
\]
a contradiction for sufficiently large $q$.
\end{proof}

\begin{proof}[Proof of Theorem~\ref{thm:main}]
Colour the edges of $G_1$ with the first colour, the edges of
$G_2\setminus G_1$ with the second colour, and all remaining edges of
the complete graph with the third colour. By
Proposition~\ref{prop:construction}, this colouring gives
$r(4,4,k)>N$.

For sufficiently large $t$, Bertrand's postulate provides a prime $q$
with
\[
 \frac{t}{4(\log t)^{1+\eta}}
 \le q\le
 \frac{t}{2(\log t)^{1+\eta}}.
\]
Then $k\le t$ and $\log q\sim\log t$, so
\[
 r(4,4,t)\ge c_\eta\frac{t^5}{(\log t)^{4+5\eta}}.
\]
For $0<\eps\le5$, take $\eta=\eps/5$. The case $\eps>5$ follows by
weakening the bound.

For the upper bound, apply the upper-bound part of
\cite[Theorem~4 and Appendix~B]{HW} with
$S=(4-2)+(4-2)=4$. This gives
\[
 r(4,4,t)\le C\frac{t^5}{(\log t)^4}.
\]
As noted in the introduction, this upper bound does not use the
pseudorandom graph existence assumption in their lower-bound result.
Combining it with the lower bound for every fixed $\eps>0$ proves the
asserted logarithmic exponent $4+o(1)$.
\end{proof}

\appendix

\section{The single-colour source interfaces}\label{app:source}

All references to the OpenAI source are pinned to the commit
\begin{center}
\src{adc7f1241b42e322a6451854ab7e4b4c146bf78a}.
\end{center}
This appendix records the interface substitutions; it does not reprove
the scanning argument, geometric classification, or window compression.
Take $P\ge\max\{64,\src{stepPrefactor(3)}\}$, which is the fixed lower
bound required in the main text.

Identify a flag $(p,H)$ with the point--dual-point pair
$(p,[\ell_H])$, where $\ker\ell_H=H$. This identification is injective
and preserves entropy and candidate-set cardinalities. The flag
condition supplies the incidence assumption in the original interface,
and \eqref{eq:legal} supplies scan consistency.

\subsection{Nonnegativity and the signed upper bound}

Set $J_{\mathrm{prev}}=3\sigma+\log P+\Delta$ and
$J_0=3\sigma+\log64$. The current table entries satisfy the original
assumption $\log|S(u,j)|\le J_{\mathrm{prev}}$. In~\cite{OAI-scan},
use a four-dimensional vector space: the dimension-related natural
number parameter in \src{twoScan_subset_deficit} is set to $1$, while
$n$ continues to denote the sequence length in this paper. Set all the
selected-position sets in that theorem to $\varnothing$. The theorem
then gives
\begin{equation}\label{eq:excess-nonnegative}
 0\le\mathfrak e,\qquad
 \mathfrak e=\src{excessCost}(F,u,J_{\mathrm{prev}},J_0).
\end{equation}
Only the \emph{test position sets} in the theorem are made empty; the
original sequence is not deleted. The empty set is disjoint from the
expensive positions, and the entropy of the empty subsequence is zero.
Thus the left-hand side of the source inequality is zero, without any
length-packing or block-occupancy condition.

By the definition in the same file and \src{twoMask_log_card}
in~\cite{OAI-budget},
\begin{equation}\label{eq:source-excess}
 \mathfrak e=\HH(u)+2n\log15+nJ_0-\HH(F)
             +\mathbb E r(F)(J_{\mathrm{prev}}-J_0).
\end{equation}
Here $r(F)$ counts the positions that require additional recording in
the two-directional scan in the source. We call these the
\emph{expensive positions}; their selection need not be reproduced in
the main text. The bound \src{twoExpensive_linear}
in~\cite{OAI-budget}, with dimension parameter $d=1$, together with
$r(F)\le n$, gives
\[
 0\le\mathbb E r(F)\le\min\{n,8000q\sigma\}.
\]

Split $J_{\mathrm{prev}}-J_0=\Delta+\log(P/64)$. Use the bound
$8000q\sigma$ for the coefficient of $\Delta$ and the bound $n$ for the
constant term. Both coefficients are nonnegative, since $P\ge64$ and
$\Delta\ge0$. Hence
\begin{align*}
 \mathfrak e
 &\le\log|\CC|+3n\sigma-\HH(F)
       +(2\log15+\log P)n+8000q\sigma\Delta\\
 &\le\DD(F,\CC,\Delta).
\end{align*}
Together with \eqref{eq:excess-nonnegative}, this proves
Proposition~\ref{input:one}\textup{(i)}. Equivalently, apply
\src{excessCost_linear} with
$A=(3n\sigma-\HH(F))/n$. Its statement does not assume $A\ge0$.
Thus allowing a signed entropy difference does not strengthen the
original interface.

\subsection{Compression parameters and finite table families}

In \src{eventually_selected_step} from~\cite{OAI-step}, set $i=0$;
the number of classes is then $24$. Use $n$ as the original input
length and $m=\lfloor n/48\rfloor$ as the auxiliary length. Since
$n\ge aQ$, sufficiently large $q$ satisfies
\[
 n/96\le m\le n/48,\qquad 24m+8000q\sigma\le n.
\]
Indeed, $24m\le n/2$ and $8000q\sigma=o(n)$, so the length-packing
condition holds. Take the fixed length constants in the source theorem
to be $a/96$ and $1$.

Since $\mathfrak e\le\DD$, the budget \eqref{eq:input-budget} implies,
with $C_0=24\cdot96B_0$, that
\[
 \mathfrak e\le C_0mD\sigma^{-\beta}/24.
\]
The block-occupancy bound is also covered because $C_0\ge B_0$. Choose
the even integer $R=2\lceil\sigma^\beta/2\rceil$. This meets the
condition \src{Admissible} in~\cite{OAI-scales}, with
$\src{scaleKstar}=D\sigma^{6\beta}$. Legality and the block conditions
correspond to the source assumptions through
\src{rectangle_pair_iff} and the related interfaces
in~\cite{OAI-geometry}.

The original output is an auxiliary finite probability experiment
retaining at least $n/6144$ ordered positions. Its message $M$ satisfies
$\HH(M)\le h:=D\sigma^{-\eta/3}n$, and the probability that its source
is any prescribed input configuration is at most $96/|\Omega|$.
Write $p_M$ for the message distribution. Retain only positive-probability
message values $z$ satisfying $-\log p_M(z)\le2h+2$. Since
$\mathbb E[-\log p_M(M)]=\HH(M)\le h$, Markov's inequality implies
that at least half the output probability is retained. Each retained
message has probability at least $e^{-2h-2}$, so there are at most
$e^{2h+2}$ retained message values.

Let $\Omega'$ be the image of the retained outputs in the original
configuration set. Then
\[
 \tfrac12\le\PP(\text{the output is retained})
 \le\frac{96|\Omega'|}{|\Omega|},
 \qquad |\Omega'|\ge|\Omega|/192.
\]
For every $\omega\in\Omega'$, fix one output witness, thereby fixing its
position subsequence and its new table index. This gives
\eqref{eq:output}. We then discard the auxiliary experiment; the next
step uses the uniform distribution on the original configuration
subset $\Omega'$.

The structure \src{ContextOutput} in~\cite{OAI-output} defines
\src{domains} as a function of the new message and the new position,
so the old table index need not be carried along. If the original
candidate sets consist of arbitrary point--dual-point pairs, intersect
them with the image of the flags before pulling them back. This cannot
increase their sizes. Actual labels may repeat throughout; it is the
original vertex positions that are distinct.

\subsection{Uniform thresholds and dependencies}

First fix $\eta,P,c,T,a,\delta$. The replacement procedure uses only
the fixed values $K=1,30$, together with the fixed length and
configuration proportions after the local halving step and the
corresponding constants $B_0,C_0$. Finally, choose $q$ larger than all
resulting thresholds at once. After these constants have been fixed,
\src{eventually_selected_step} is uniform over the input tables,
probability spaces, and parameters $D,J$. There is therefore no choice
of $q$ depending on a random state generated later in the argument.

The only substantial source inputs are the single-colour scanning
excess cost and the single-colour compression theorem, which lie in
the same dependency chain. No separate candidate-sequence counting
theorem or terminal entropy theorem is required.

\section{Scope of related results}\label{app:other}

This appendix records which cited results are used directly and why
several other available statements do not, without additional work,
replace the single-colour compression input.

\subsection{Inputs used in the main argument}

Brada\v{c}'s result~\cite[Lemma~2.11]{Bradac} directly supplies the
label digraph without a transitive four-vertex tournament.
He and Wigderson~\cite[Appendix~B]{HW} directly supply the upper bound.
The nonnegativity of the single-colour scanning cost is taken
from~\cite{OAI-scan}; it is not obtained indirectly through a
prefix-tree counting argument. Subadditivity of entropy is used
from~\cite[Theorem~2.6.6]{CT}, without reproving it. The two-colour
argument consists of \eqref{eq:balance}, the local replacement step,
and the finite-round cost calculation.

\subsection{Related container theorems}

\paragraph{Campos and Samotij.}
Theorem~B of~\cite{CS} requires no codegree assumptions. For every
hypergraph and $0<p\le\delta<1$, it supplies a fingerprint $S$ and a
container $C$ such that $|S|\le p|V|/\delta$ and
\[
 \PP(S\cup C_p\text{ is independent})
 \ge(1-p)^{\delta|C\setminus S|}.
\]
It does not directly assert that $|C|$ is small. To use this statement
instead of compression in the present setting, one would still need
to prove that the probability lower bound forces small candidate sets
at each position and to account for the number of fingerprints.
Theorem~D in the same paper expresses these requirements in terms of
the absence of a suitable balanced supersaturation subhypergraph;
it does not verify them automatically. Thus the absence of codegree
assumptions in Theorem~B alone does not remove the need for the
single-colour compression input here.

\paragraph{Balogh and Samotij; Nenadov and Pham.}
Theorem~1.6 of Balogh and Samotij~\cite{BS} packages a multiround
container procedure, but still requires codegree bounds at every
order. To turn its output of a large subset with few edges into a
small container, one needs a supersaturation statement at the target
scale. Theorem~1.1 of Nenadov and Pham~\cite{NP} gives a short proof
of the container theorem, but its input still requires a
$(p,K)$-uniformly-spread measure supported on hyperedges: for every
nonempty $L$,
\[
 \nu(\{E:L\subseteq E\})\le\frac{Kp^{|L|-1}}{|V|}.
\]
A low-entropy candidate table in the present proof is not such a
measure on hyperedges, so the theorem cannot be applied by a direct
parameter substitution. These results may be useful in a different
proof, but they are not interchangeable inputs for the present one.

\paragraph{Nenadov.}
Theorem~1.1 of~\cite{N} is the closest of these statements to a
container result that treats only most inputs. It permits an
exceptional family of relative size $\eps^m$ and assumes the
second-moment condition
\[
 \Delta_{\mathcal H}(m_0)
 \le B\mu_{\mathcal H}(m_0)^2/m_0.
\]
However, the legal sequences selected in our argument are not uniform
samples from the family of all $m$-element sets, and could be
concentrated entirely in the exceptional family. In the most direct
independent-tuple count, the required relative probability is about
$q^{-2m}$, for which an error of $\eps^m$ with fixed $\eps$ is too
large. Allowing $\eps$ to shrink with $q$ would require new quantitative
control of the threshold $T(\eps)m_0$. The paper itself distinguishes
this probabilistic version from full-cover container results suited
to iteration.

\subsection{Independent-set counts and the target scale}

For graphs containing no induced $bK_a$, Buci\'{c}, Chudnovsky and
Codsi~\cite[Theorem~1]{BCC} prove a bound of the form
\[
 i(G)\le |V(G)|^{O(1)}\alpha(G)^{O(\alpha(G))},
\]
where $i(G)$ denotes the number of independent sets of $G$. Taking
$a=4$ and $b=1$ includes all $K_4$-free graphs. However, $\alpha(G)$ in
this estimate is the \emph{actual} independence number of the entire
graph, not the target value $t$ that we wish to exclude. One cannot
substitute $\alpha(G)=t$ without a proof, nor does the total number of
independent sets directly provide the required counts for ordered,
possibly repeated labels and position-wise candidate sets.

The full tuple-counting estimate in Brada\v{c}'s
Lemma~2.12~\cite{Bradac} still requires $m\ge Cq(\log q)^2$. It does
not reach the scale $q(\log q)^{1+\eta}$ for arbitrarily small fixed
$\eta>0$. Claim~2.13 of that paper remains available for counting with
restricted candidate sets, but the present proof does not use that
route.

\paragraph{Scope of this comparison.}
None of the statements discussed above directly replaces the complete
candidate-table iteration in the present argument. We therefore use
the single-colour scanning and compression interfaces explicitly,
without treating unverified supersaturation, codegree, or
exceptional-family estimates as implicit assumptions.

\begin{thebibliography}{99}
\small\raggedright
\setlength{\itemsep}{4pt}

\bibitem{BS}
J.~Balogh and W.~Samotij.
An efficient container lemma.
\emph{Discrete Analysis}, 2020, Paper No.~17.
\href{https://arxiv.org/abs/1910.09208}{arXiv:1910.09208}.
Theorem~1.6.

\bibitem{Bradac}
D.~Brada\v{c}.
Off-diagonal Ramsey numbers.
\href{https://arxiv.org/abs/2605.28793v3}{arXiv:2605.28793v3},
16 June 2026.
Lemma~2.11 is used in the main argument; Lemmas~2.7 and~2.12 and
Claim~2.13 are relevant for comparison.

\bibitem{BCC}
M.~Buci\'{c}, M.~Chudnovsky, and J.~Codsi.
Counting independent sets in structured graphs.
\href{https://arxiv.org/abs/2406.07799v1}{arXiv:2406.07799v1}, 2024.
Theorems~1 and~2.

\bibitem{CS}
M.~Campos and W.~Samotij.
Towards an optimal hypergraph container lemma.
\href{https://arxiv.org/abs/2408.06617v2}{arXiv:2408.06617v2},
19 September 2024.
Theorems~B and~D in this fixed version; the conjectural statements
are not used.

\bibitem{CT}
T.~M. Cover and J.~A. Thomas.
\emph{Elements of Information Theory}.
Wiley, 1991.
Theorem~2.6.6.
\href{https://sites.stat.columbia.edu/liam/teaching/neurostat-spr11/papers/EM/Cover\%26Thomas-Ch2.pdf}{Public scan of Chapter~2}.
Only subadditivity of entropy is used; changing the logarithm base
does not affect this inequality.

\bibitem{HW}
X.~He and Y.~Wigderson.
Multicolor Ramsey numbers via pseudorandom graphs.
\href{https://arxiv.org/abs/1910.06287v3}{arXiv:1910.06287v3}.
The upper-bound part of Theorem~4 and Appendix~B.

\bibitem{N}
R.~Nenadov.
Probabilistic hypergraph containers.
\emph{Israel Journal of Mathematics}, \textbf{261} (2024), 879--897.
\href{https://arxiv.org/abs/2111.06363v2}{arXiv:2111.06363v2}.
Theorem~1.1.

\bibitem{NP}
R.~Nenadov and H.~T. Pham.
Short proof of the hypergraph container theorem.
\href{https://arxiv.org/abs/2408.08514v3}{arXiv:2408.08514v3},
10 September 2024.
Theorem~1.1 and Lemma~2.1.

\bibitem{OAI-scan}
OpenAI.
\emph{SharpLogRamsey source}, fixed commit specified in
Appendix~\ref{app:source}.
\oaifile{Entropy/MarkingTwoEntropy.lean}{Entropy/MarkingTwoEntropy.lean}:
\src{excessCost} and \src{twoScan_subset_deficit}.

\bibitem{OAI-budget}
OpenAI.
\emph{SharpLogRamsey source}, the same fixed commit.
\oaifile{Marking/ExcessCost.lean}{Marking/ExcessCost.lean}:
\src{twoMask_log_card}, \src{twoExpensive_linear}, and
\src{excessCost_linear}.

\bibitem{OAI-step}
OpenAI.
\emph{SharpLogRamsey source}, the same fixed commit.
\oaifile{Marking/SelectedStep.lean}{Marking/SelectedStep.lean}:
\src{eventually_selected_step}.
\oaifile{Marking/ClassStep.lean}{Marking/ClassStep.lean}:
\src{stepPrefactor}.

\bibitem{OAI-output}
OpenAI.
\emph{SharpLogRamsey source}, the same fixed commit.
\oaifile{Execution/ContextOutput.lean}{Execution/ContextOutput.lean}:
\src{ContextOutput}, in particular \src{source}, \src{chosen},
\src{domains}, \src{dominated}, and \src{cost}.

\bibitem{OAI-scales}
OpenAI.
\emph{SharpLogRamsey source}, the same fixed commit.
\oaifile{Parameters/SourceScales.lean}{Parameters/SourceScales.lean}:
\src{beta}, \src{Admissible}, and \src{scaleKstar}.

\bibitem{OAI-geometry}
OpenAI.
\emph{SharpLogRamsey source}, the same fixed commit.
\oaifile{Construction/InitialRectangles.lean}{Construction/InitialRectangles.lean};
\oaifile{Construction/InitialTools.lean}{Construction/InitialTools.lean};
\oaifile{Marking/OrderedClassLaw.lean}{Marking/OrderedClassLaw.lean}.
The block counts, \src{rectangle_pair_iff}, and
\src{scanConsistent_ordered} are used.

\end{thebibliography}

\end{document}
